\documentclass[abstracton,parskip=half]{scrreprt}
\usepackage[english,ngerman]{babel}
\usepackage[utf8]{inputenc}
\usepackage[T1]{fontenc}
\usepackage{graphicx}
\usepackage{hyperref}
\usepackage{todonotes}
\usepackage{listings}
\usepackage{url}
\usepackage{amssymb}




\lstdefinelanguage{scala}{
  keywords={%
          abstract,case,catch,class,def,do,else,extends,%
          false,final,finally,for,forSome,if,implicit,import,lazy,%
          match,new,null,object,override,package,private,protected,%
          return,sealed,super,this,throw,trait,true,try,type,%
          val,var,while,with,yield},
  otherkeywords={=>,<-,<\%,<:,>:,\#,@},
  sensitive=true,
  morecomment=[l]{//},
  morecomment=[n]{/*}{*/},
  morestring=[b]",
  morestring=[b]',
  morestring=[b]"""
}[keywords,comments,strings]


\lstset{
      basicstyle=\small,%
      frame=tb,
      basewidth={0.5em,0.45em},
      literate={=}{{$=$}}1
               {>}{{$>$}}1 {<}{{$<$}}1 {\\}{{$\lambda$}}1
               {\\\\}{{\char`\\\char`\\}}1
               {->}{{$\rightarrow$}}2 {>=}{{$\geq$}}2 {←}{{$\leftarrow$}}2
               {<=}{{$\leq$}}2 {⇒}{{$\Rightarrow$}}2
               {|}{{$\mid$}}1
}

\hypersetup{
   hidelinks=true,
   pdfauthor={Julian Andres Klode},
   pdftitle={Bachelor Thesis: Performance and Interfaces of Datatype-Generic Scala Programs},
   pdfsubject={Programming Languages},
}

\author{Julian Andres Klode\\Matriculation number: 2403668\\\\Supervisor: Yufei Cai}
% Let's abuse the subject field for now
\subject{\includegraphics[height=.13\textheight]{graphics/Logo_uni}\\Programming Languages and Software Technology Group\\Prof. Ostermann\\[1.5cm]}
\title{Performance and Interfaces of Datatype-Generic Scala Programs}
\subtitle{Bachelor Thesis}
\newcommand{\cd}{\texttt}
\newcommand{\ex}{\paragraph{Example:}}
\date{}


\begin{document}
    \pagenumbering{roman}
    \maketitle{}
    \begingroup
\let\titlepage\par
\let\endtitlepage
\let
\selectlanguage{ngerman}
\begin{abstract}
Generische Programmierung ist eine Vorgehensweise, die es ermöglicht, Funktionen
über die Struktur eines Datentypes zu definieren, anstelle von eines spezifischen
Datentypes. Dies erlaubt es einer Funktion auf beliebigen Datentypen zu operieren,
selbst auf solchen, die dem Programmier der Funktion noch nicht bekannt waren.

In dieser Bachelorarbeit werde ich Portierungen von drei Haskell Bibliotheken
(LIGD, EMGM, Uniplate) nach Scala vorstellen und ihre Funktionalität und Geschwindigkeit miteinander
sowie mit Shapeless vergleichen.

Ich werde außerdem zeigen, dass LIGD in Scala erweitert werden kann, sodass sie
beinahe so mächtig wie EMGM ist, indem ich das Problem der Erweiterbarkeit
von Funktionen löse, welches bisherige Implementierungen stark beeinträchtige.
\end{abstract}
\selectlanguage{english}
\begin{abstract}
Generic programming is a technique that allows functions to be defined over
the shape of an type, rather than a specific type; thus allowing one function
to be used on objects of arbitrary types, even those that were not known to
the programmer of the function.

In this thesis, I will present ports of three Haskell libraries (LIGD, EMGM, Uniplate) to Scala,
and compare their features and performance against each other and
against shapeless, a native Scala library.

I will show that LIGD in Scala can be extended to be almost as powerful as EMGM;
solving the issue with extensibility that made it vastly inferior in previous
implementations.


\end{abstract}
\endgroup

    \selectlanguage{english}

    \tableofcontents{}
    \clearpage
    \pagenumbering{arabic}
    \setcounter{page}{1}

    \chapter{Introduction}

Many\footnote{Such as Haskell, Java, Scala, C++} programming languages these days support a form of parametric polymorphism.
Parametric Polymorphism allows the definition of data types that abstract over
the type of some element they store, such as lists that can store elements of
some unspecified type \lstinline{a}:
\begin{lstlisting}[language=scala]
sealed abstract class List[a]
case class Nil[a]() extends List[a]
case class Cons[a](car: a, cdr: List[a]) extends List[a]
\end{lstlisting}
In Scala and other object-oriented (or influenced) languages this concept is
called generics.
Datatype-generic programming is a different concept. Instead of abstracting
over the type of the elements, it abstracts over the container type.

For example, consider we have types to represent arrays and lists. Consider that
those store integers. We want to increase every integer by one. That is, we want
to write a function that takes a container of integers and increases every integer
by one. That is, we want to write something like:

\begin{lstlisting}[language=Scala]
def incInt[C[_]](c: C[Int]) = everywhere((i: Int) => i + 1)(c)
\end{lstlisting}
(Note: This does not actually work with shapeless)

Datatype-generic programming does not only abstract containers. In Haskell and
other functional programming languages, data types (except some built-in scalar
types like integers and floating point numbers) are represented using
algebraic data types, that is (labelled) sum types and product types.
In Scala, case classes with a sealed abstract base class (or trait) can
represent algebraic data types.\cite{scalagp}
But Scala also allows unsealed classes, so further cases can be added even in
different modules.


\section{Existing Haskell Modules}
\subsection{\texttt{Data.Generics}}
\subsection{\texttt{Unbound}}

\cite{DBLP:conf/tldi/LammelJ03}



    \chapter{Porting from Haskell to Scala}

    This chapter introduces some basic concepts for translating Haskell code
    to Scala. These concepts were explained by Oliveira and Gibbons~\cite{scalagp}, this
    chapter summarizes them.

    \section{Type classes and Instances}

    Haskell as a concept of type classes and instances, whereas Scala offers
    traits. Traits differ from type classes in that their methods always contain
    an implicit \cd{this} object, that is, they combine data and functions;
    whereas type classes do not reference any data.

    Given a type class like
    \begin{lstlisting}[language=Haskell,gobble=6]
      class MyEq a where
        myEquals :: a -> a -> Boolean
        myEquals a b = False
    \end{lstlisting}
    how do we translate this to Scala? There are two approaches:

    \subsection{Object-oriented translation}
    The object-oriented approach uses traits just like interfaces in
    Java are used. For example, the \cd{MyEq} type class would be translated
    to the following trait:
    \begin{lstlisting}[language=Scala,gobble=6]
      trait MyEq[T] {
        def myEquals(b: T) : Boolean = false
      }
    \end{lstlisting}
    In order to implement this trait for a class the class must extend the
    trait and override the methods of the trait:
    \begin{lstlisting}[language=Scala,gobble=6]
      class MyInt(val i: Int) extends MyEq[MyInt] {
        override def myEquals(b: MyInt) : Boolean = i == b.i
      }
    \end{lstlisting}
    This already shows one deficiency of the traits approach: We cannot
    implement a trait for a type without modifying that type itself.

    Writing a function using this approach is straight forward:
    \begin{lstlisting}[language=Scala,gobble=6]
      def myEquals[T <: MyEq[T]](a: T, b: T) = a.myEquals(b)
    \end{lstlisting}

    This style makes it easy to specify contexts, so for example, if we have
    a class \cd{MyOrd} for ordering that also requires equality, we can write
    \begin{lstlisting}[language=Scala,gobble=6]
      trait MyOrd[T] extends MyEq[T] {
        def myLessOrEqual(b: T) : Boolean
      }
    \end{lstlisting}

    If we want to translate a type class containing a method that does not take a
    parameter of the type A, this style of writing becomes inconsistent.

    \ex Consider the type class
    \begin{lstlisting}[language=Haskell,gobble=6]
      class PerformerFactory a where
        factor :: a
        perform :: a -> Boolean
    \end{lstlisting}
    We can only directly translate the \cd{perform} method to Scala. The \cd{factor}
    method can only be translated by use of a generic function -- this will
    require a stub implementation, that instances can specialize.
    \begin{lstlisting}[language=Scala,gobble=6]
      trait PerformerFactory[A] {
        def factor : A
      }
      def perform[A](a: A) : Boolean = throw new Exception
    \end{lstlisting}
    The end result is inconsistent, because one method is now a method of
    instance objects whereas the other is translated to a global function. It
    also permits instances of the trait that only implement the \cd{factor}
    method.
    \subsection{The functional approach}
    Remember the type class MyEq:
    \begin{lstlisting}[language=Haskell,gobble=6]
      class MyEq a where
        myEquals :: a -> a -> Boolean
        myEquals a b = False
    \end{lstlisting}
    We can ignore the implicit this parameter in traits and translate the type
    class as:
    \begin{lstlisting}[language=Scala,gobble=6]
      trait MyEq[A] {
        def myEquals(a: A, b: A) = false
      }
      def myEquals[T](a: T, b: T)(implicit eq: MyEq[T]) = eq.myEquals(a, b)
    \end{lstlisting}
    We also provided a global \cd{myEquals} function to make our life easier.

    Now in order to implement the trait for a type, we need to define an instance
    of the trait for this type. Instances of those traits correspond to dictionaries
    in Haskell, and can be implemented like this:
    \begin{lstlisting}[language=Scala,gobble=6]
      implicit def MyEqInt = new MyEq[Int] {
        override def myEquals(a: Int, b: Int) = a == b
      }
    \end{lstlisting}
    The \cd{implicit} keyword is a special syntactic sugar provided by Scala
    that makes our work easier. Without it, if we wanted to compare two objects,
    we would need to manually pass the instance of the trait to the function,
    like this:
    \begin{lstlisting}[language=Scala,gobble=6]
      > myEquals(1, 1)(MyEq)
      true
    \end{lstlisting}
    By making \cd{MyEqInt} and the \cd{eq} parameter of \cd{myEquals} implicit,
    the Scala compiler can automatically figure out that \cd{MyEqInt} is the
    correct argument for the \cd{eq} parameter, and we can now run:
    \begin{lstlisting}[language=Scala,gobble=6]
      > myEquals(1, 1)
      true
    \end{lstlisting}


    We can also implement instance that require some sort of context. For
    example, if we want to implement \cd{MyEq} for tuples of objects implementing
    \cd{MyEq} as well, we can write:
    \begin{lstlisting}[language=Scala,gobble=6]
      implicit def myEqTuple[A,B](implicit ea: MyEq[A], eb: MyEq[B]) =
        new MyEq[(A,B)]{
        override def myEquals(a: (A,B), b: (A,B)) = ea.myEquals(a._1, b._1)
                                        && eb.myEquals(a._2, b._2)
      }
    \end{lstlisting}
    Now we can compare tuples: For example, in the following example, Scala will
    automatically infer the correct instance of \cd{MyEq[(Int, Int)]}:
    \begin{lstlisting}[language=Scala,gobble=6]
      > myEquals((1,2), (1,2))
      true
    \end{lstlisting}
    We can still pass  it manually though
    \begin{lstlisting}[language=Scala,gobble=6]
      > myEquals((1,2), (1,2))(myEqTuple(MyEqInt, MyEqInt))
      true
    \end{lstlisting}

    \section{Making the most of type inference}

    Compared to Haskell, Scala's type inference is weaker. For example, when
    defining a function, the type of parameters need to be defined. An exception
    to this are anonymous functions. If the Scala compiler knows which type the
    function must have, it can automatically infer the type of an anonymous
    function. Thus, in order to avoid repetitious type information in the code,
    we can change our \cd{MyEq} example to:
    \begin{lstlisting}[language=Scala,gobble=6]
      trait MyEq[A] {
        def myEquals : (A => A => Boolean) = a => b => false
      }
      implicit def MyEqInt = new MyEq[Int] {
        override def myEquals = a => b => a == b
      }
    \end{lstlisting}
    Type annotations are not required in method implementations anymore. This
    technique is shown in the port of the EMGM library.

    \chapter{Libraries}

    \section{LIGD -- Translation of LIGD}
    \lstinputlisting[language=scala,firstline=55,lastline=70,float=ht,caption=LIGD Basics,label=ligd-base]{scala/ligd.scala}

    LIGD is a simple library that provides a \cd{Rep} trait that describes
    the structure of a type using sums, products, and scalar values.

    Listing~\ref{ligd-base} shows the basic \cd{Rep} trait with instances for
    scalar types, sums (represented using \cd{Either}), and products (represented
    using \cd{Tuple2}). As explained in the previous chapter, we use implicit
    here to allow the compiler to automatically infer the correct \cd{Rep}
    instance for generic function calls.



    \subsection{User-defined types}
    In order to work with user-defined (case) classes, those need to be
    translated to a representation using sums and products. We need an
    isomorphism between the user-defined class and a sum of product type. First
    of all, we can define an isomorphismus as:
    \begin{lstlisting}[language=scala,gobble=6]
      class EP[B, C](val from: B ⇒ C, val to: C ⇒ B)
    \end{lstlisting}
    Now we can define a \cd{Rep} instance, that represents values of a custom
    type \cd{B} and converts it to a sum of product notation of type \cd{C},
    something like this:
    \begin{lstlisting}[language=scala,gobble=6]
      case class RType[C, B](val c: Rep[C], val ep: EP[B, C]) extends Rep[B]
    \end{lstlisting}
    There is one problem though: The representation of \cd{C} may be recursive,
    for example a list type is a sum of nothing, and a product of an element and
    the list type. Thus, the parameter \cd{c} must be lazy. Scala does not allow
    lazy \cd{val}s on case classes, though, so this needs to be worked around.

    There are two approaches: One is encoding the thunk manually, like:
    \begin{lstlisting}[language=scala,gobble=6]
      case class RType[C, B](val c: () => Rep[C], val ep: EP[B, C]) extends Rep[B]
    \end{lstlisting}
    This is very easy, but looks inconsistent in usage. Another approach is to
    not use a case class and provide a companion object with custom \cd{apply}
    and \cd{unapply} methods. Listing~\ref{ligd-rtype} shows how this is
    implemented.
    \lstinputlisting[language=scala,firstline=78,lastline=97,float=ht,caption=Implementation of RType,label=ligd-rtype]{scala/ligd.scala}

    \textbf{Example (Lists):} In order to represent a list, we encode the
    list as a sum of products, or to be precise, we want to encode a \cd{List[A]}
    as: \cd{Either[Unit, (A, List[A])]}. In order to do this, we need functions
    that convert from \cd{List[A]} to \cd{Either[Unit, (A, List[A])]} and
    vice versa. An implementation of those functions is given in figure~\ref{ligd-rlist}.
    %-- We can then define an RType instance as:
    %\lstinputlisting[language=scala,firstline=109,lastline=113]{scala/ligd.scala}

    \lstinputlisting[language=scala,firstline=103,lastline=124,float=ht,caption=Implementation of rList,label=ligd-rlist]{scala/ligd.scala}

    Note how the first parameter of \cd{RType} is recursive.

    \subsection{Writing generic functions}
    Writing generic functions using LIGD is simple. We just pass the object(s)
    we are interested in to the function, along with their type representations,
    and then pattern match on the type representations.

      \lstinputlisting[language=scala,firstline=37,lastline=52,float=ht,caption=Generic Equality in LIGD,label=ligd-geq]{scala/ligd.scala}

    This function is pretty much self-explaining.

    \subsection{Differences to the Haskell implementation}
    One minor difference is that we implemented \cd{Rep} like we would implement
    a type class in Haskell. The original paper manually passed around \cd{Rep}
    instances for generic functions, we make use of Scala's implicit values.


    \subsubsection{Ad-hoc cases using subtyping}
    A key difference that extends the functionality of LIGD in Scala is that
    Scala permits sub-typing. If we look at our definition of RType, we see
    that it is not \cd{sealed}, allowing it to be extended by subclasses.

    We can use this to create custom sub classes for types we are interested
    in, and pattern match on them. Because those are sub classes of \cd{RType},
    existing code will be unaffected by new types.

    An example for this is the list type. Because we extended \cd{RType} as
    \cd{RList} we can now write a generic function that finds all lists of a
    specific element type in a data structure.

    \subsubsection{Equality comparisons of type representations}
    The Haskell version of LIGD does not support comparing Rep instances
    against each other. This is quite useful, though, as it allows to
    implement simple variants of generic operations that look for objects
    of a specific type in another object.


    When representing generic types such as Lists, we can create a subclass
    of RType with parameters representing the generic arguments. We can then
    define a custom \cd{equals()} method, as seen in listing~\ref{rmycustomtype}.

    \begin{lstlisting}[language=Scala,gobble=8,float,caption=Representing custom generic types,label=rmycustomtype]
        class RMyCustomType[T1, ...](r1: Rep[T1], ...) {
            override def equals(b: Any) : Boolean = b match {
                case RMyCustomType(r1, ...) => r1 == this.r1 &&
                                               r2 == this.r2 && ...
                case _ => false
            }
        }
    \end{lstlisting}



    \section{EMGM -- Translation of Extensible and Modular Generics for the Masses}

    \section{Uniplate}
    Uniplate is a library that differentiates between two basic type of
    objects: Uniplates and Biplates. Uniplates are simple recursive types
    like terms or expressions, which form a graph where all inner nodes have
    the same common supertype. Biplates are types that contain other Uniplates,
    for example, a \cd{List a} has a \cd{Biplate List a}.


    Uniplate can theoretically be ported from Haskell to Scala, but it is not
    feasible to do without a way to automatically generate instances of the
    \cd{Biplate} type class, because the amount of biplates needed is quadratic
    to the number of types.

    A basic implementation of uniplate can be found in the \cd{uniplate.scala}
    source file in section~\ref{src:uniplate}.


    \section{Shapeless -- Scala Implementation of Scrap your Boilerplate}

    Shapeless is an adaption of Scrap your Boilerplate\cite{DBLP:conf/tldi/LammelJ03}
    from Haskell to Scala. Compared to the Haskell implementation, the Scala
    implementation simplifies some things a bit.

    Comparing our ports of other libraries against shapeless helps us determine
    whether the solutions are efficient.


    \chapter{Implementing Datatype-Generic Programming in Scala}
    \section{Name-binding Operations}
    \section{Derived Data Types}
    \section{GPBench}
    GPBench is a suite of benchmarks for comparing generic programming
    libraries in Haskell \cite{DBLP:conf/haskell/RodriguezJJGKO08}.

    \section{Benchmarks}

    To test the performance of the libraries, we ran various tests with them. For
    each test, we executed it 600 thousand times, and calculated the average
    time of the last 300 thousand executions.

    \begin{table}[ht]
    \begin{tabular}{c|r|r|r|r}
\bf{test}  & \bf{company}    & \bf{geq}        & \bf{min}        & \bf{sum}       \\
\hline
Direct     &          3,0 μs &          0,4 μs &          0,6 μs &          0,4 μs\\
Shapeless  &         39,3 μs &             N/A &          4,7 μs &          4,8 μs\\
LIGD       &         16,0 μs &         15,3 μs &         13,5 μs &         13,8 μs\\
EMGM       &         18,4 μs &         16,5 μs &         12,2 μs &         16,5 μs\\
\end{tabular}

    \caption{Benchmark results}
    \label{bench}
    \end{table}

    As we can see in in table~\ref{bench}, all tested libraries perform roughly
    the same time-wise.


    \chapter{Conclusion}

    \begin{table}[t!]
      \begin{tabular}{c|cccc}
        Feature (and test)    & LIGD & EMGM & Uniplate & Shapeless \\
                     \hline
        Ad-hoc cases (company) & \checkmark & \checkmark & \checkmark & \checkmark \\
        Multi-parameter functions (geq) & \checkmark & \checkmark & ($\times$) & \checkmark
      \end{tabular}

      \caption{Feature overview}
    \end{table}

    \bibliographystyle{alpha}
    \bibliography{thesis}

    \appendix{}
    \chapter{Source Code}

    \lstset{frame=none}
    \section{LIGD}
    \lstinputlisting[language=scala,firstline=19,label=src:ligd,showstringspaces=false]{scala/ligd.scala}
    \section{EMGM}
    \lstinputlisting[language=scala,firstline=19,label=src:emgm,showstringspaces=false]{scala/emgm.scala}
    \section{Uniplate}
    \lstinputlisting[language=scala,firstline=19,label=src:uniplate,showstringspaces=false]{scala/uniplate.scala}
    \section{Test Data}
    \lstinputlisting[language=scala,firstline=19,label=src:data,showstringspaces=false]{scala/data.scala}
    \section{Company tests}
    \lstinputlisting[language=scala,firstline=19,label=src:companytest,showstringspaces=false]{scala/CompanyTest.scala}

    \section{Benchmark}
    \lstinputlisting[language=scala,firstline=19,label=src:test-,showstringspaces=false]{scala/test.scala}
\end{document}
